
%(BEGIN_QUESTION)
% Copyright 2006, Tony R. Kuphaldt, released under the Creative Commons Attribution License (v 1.0)
% This means you may do almost anything with this work of mine, so long as you give me proper credit

En konge får en skinnende krone i gave.  Giveren påstår at den er av rent, massivt gull.  Den ser ut som gull, men kongen -- klok som han er -- vet at den kan være gullbelagt tinn eller et annet billig metall.

Han gir kronen til en berømt vitenskapsmann for analyse, med krav om at kronen ikke må skades av testingen.  Vitenskapsmannen vurderer hvordan sammensetningen kan bestemmes ikke-destruktivt, og bestemmer at tetthetsmåling ved veiing både tørr og neddykket i vann er tilstrekkelig, siden gull er vesentlig tyngre enn billige metaller.

Kronen veier 5 pund tørr.  Når den er helt neddykket i vann, veier den 4.444 pund.  Er den virkelig av massivt gull?

\underbar{file i00273}
%(END_QUESTION)





%(BEGIN_ANSWER)

Finn først kronens spesifikke vekt.  Hvis den veier 5 pund tørr og 4.444 pund neddykket, må vekten av fortrengt vann være 0.556 pund.  En krone som veier 5 pund, men har volum tilsvarende 0.556 pund vann, må ha tetthet omtrent 9 ganger vann (5 / 0.556 $\approx$ 9).  Vi kan si spesifikk vekt 9, eller tetthet 9 g/cm$^{3}$.

Rent gull har spesifikk vekt nær 19, så denne kronen kan ikke være av rent, massivt gull.  Den er mest sannsynlig laget av kobber (D = 8.96 g/cm$^{3}$) eller tinn (D = 7.31 g/cm$^{3}$), belagt med gull.

\vskip 10pt

Ifølge legenden er dette hvordan Arkimedes kom frem til prinsippet om sammenhengen mellom fortrengning og oppdriftskraft.

%(END_ANSWER)





%(BEGIN_NOTES)


%INDEX% Physics, static fluids: Archimedes' principle

%(END_NOTES)
